\section{Implementación experimental}

La implementación física no usó el código generado de \texttt{Controller\_BoP}. Por tiempo, el
equipo decidió probar el lazo con un firmware escrito a mano
(\texttt{firmware/bop\_controlador\_estable/}) que implementa la misma ley \eqref{eq:ley} con
las mismas ganancias. La función de la ley del primer firmware de prueba se comprobó,
compilada en la PC, contra la del MIL, con una diferencia máxima de \SI{6e-7}{\degree}, que
corresponde a la precisión simple del firmware. Sobre esa función se agregaron protecciones durante las pruebas: zona de tolerancia de
\SI{5}{\milli\metre}, tope del integrador en \SI{3}{\degree}, descarga del integrador lejos de la
referencia, límite de velocidad de \SI{6}{\degree} por cuadro, filtro de saltos de posición y un
\emph{trim} de nivel. Por eso los resultados de esta sección no son una validación del código
verificado en SIL y PIL, sino de la misma ley con esas protecciones.

\subsection{Arquitectura del lazo físico}

La cadena completa es: captura de la cámara a \SI{30}{fps} en la PC con Linux $\rightarrow$
homografía de píxeles a milímetros sobre el cuadrado de \SI{290}{\milli\metre} de las marcas del
plato $\rightarrow$ segmentación HSV de la bola con filtro de área y circularidad $\rightarrow$
corrección de paralaje $\rightarrow$ filtro de Kalman de velocidad constante $\rightarrow$ trama
binaria de \SI{20}{bytes} por el puerto serie a \SI{921600}{baudios} $\rightarrow$ controlador
en la ESP32 $\rightarrow$ saturación $\rightarrow$ conversión de inclinación a comando del servo
con la tabla medida $\rightarrow$ PWM de \SI{50}{\hertz} $\rightarrow$ servos. El
algoritmo de control se ejecuta en la ESP32; la PC solo mide y filtra. La ESP32 lleva el plato a
horizontal si la visión indica que no hay bola o si deja de recibir tramas durante
\SI{200}{\milli\second}, y en cada cuadro devuelve por el mismo puerto la inclinación aplicada,
que la PC guarda junto con la medición.

Las pruebas encontraron problemas de implementación que no estaban en el modelo y que se
corrigieron antes de medir: la ESP32 leía la trama byte a byte y trabajaba con datos unos
\SI{0.6}{\second} atrasados, los ejes de la cámara quedaron cruzados por el orden de los clics de
la calibración, y la primera tabla servo--plato tenía un desnivel en el neutro. Además, el
\texttt{loop()} de la ESP32 corre cada \SI{40}{\milli\second}, por lo que el controlador procesa
\num{25} de las \num{30} tramas por segundo.

Antes de energizar se comprobó: fuente de los servos separada de la ESP32, con interruptor a
mano y tierra común; primero se conecta la ESP32 y después la fuente; el firmware arranca
desarmado, con el plato horizontal, y solo actúa al presionar el botón BOOT; el comando del servo
se limita a $\pm\SI{40}{\degree}$; y el watchdog lleva el plato a horizontal si se pierden las
tramas. Durante las pruebas se agregaron al plato bordes más altos que impiden que la bola caiga.

\subsection{Resultados experimentales}

La Fig.~\ref{fig:exp} muestra los escalones de referencia de \SI{40}{\milli\metre} del ensayo
automático y una trayectoria circular de \SI{45}{\milli\metre} de radio y periodo de
\SI{4}{\second}, las dos con la pelota de ping pong y el controlador de la Sección~III. La
Tabla~\ref{tab:exp} resume las métricas de los escalones, calculadas desde el CSV con las
definiciones de la Sección~I, junto con las del MIL para escalones de \SI{50}{\milli\metre} (caso
D4). No es el mismo ensayo, porque en el prototipo no se corrió el ensayo de requisitos de
\SI{100}{\milli\metre}, pero sirve para comparar el orden de magnitud.

\begin{figure}[!t]
\centering
\begin{tikzpicture}
\begin{groupplot}[group style={group size=1 by 2, vertical sep=0.7cm}, estilo, height=3.2cm,
                  xmin=0, xmax=30]
\nextgroupplot[ylabel={$x$ (\si{\milli\metre})}, xticklabels={}, legend pos=north east,
               legend columns=2]
  \addplot[cAzul] table[col sep=comma, x=t_s, y=x_mm] {\G/exp_20261113_105834.csv};
  \addplot[cGris, densely dashed] table[col sep=comma, x=t_s, y=rx_mm] {\G/exp_20261113_105834.csv};
  \legend{Medida, Referencia}
\nextgroupplot[ylabel={$y$ (\si{\milli\metre})}, xlabel={Tiempo (\si{\second})}]
  \addplot[cNaranja] table[col sep=comma, x=t_s, y=y_mm] {\G/exp_20261113_105834.csv};
  \addplot[cGris, densely dashed] table[col sep=comma, x=t_s, y=ry_mm] {\G/exp_20261113_105834.csv};
\end{groupplot}
\end{tikzpicture}

\vspace{2mm}
\begin{tikzpicture}
\begin{axis}[estilo, width=0.5\columnwidth, height=0.5\columnwidth, axis equal image,
             xmin=-127, xmax=127, ymin=-127, ymax=127,
             xlabel={$x$ (\si{\milli\metre})}, ylabel={$y$ (\si{\milli\metre})}]
  \addplot[cAzul, line width=0.3pt, aligera=2,
           restrict expr to domain={\thisrow{modo}}{2:2}] table[col sep=comma, x=x_mm, y=y_mm] {\G/exp_20261113_105517.csv};
  \addplot[cRojo, dashed, domain=0:360, samples=73, line width=0.8pt] ({45*cos(x)}, {45*sin(x)});
\end{axis}
\end{tikzpicture}
\caption{Prueba funcional en el prototipo con el firmware de prueba. Arriba: escalones de
\SI{40}{\milli\metre} (\texttt{sesion\_20261113\_105834}). Abajo: trayectoria circular de
\SI{45}{\milli\metre} y \SI{4}{\second} durante \SI{157}{\second}, con el círculo pedido en rojo
(\texttt{sesion\_20261113\_105517}).}
\label{fig:exp}
\end{figure}

\begin{table}[!t]
\caption{Escalones en el prototipo (\texttt{sesion\_20261113\_105834}, tramos de
\SI{4.8}{\second}) y en MIL (D4). $\Delta$: cambio pedido desde la posición inicial; $e_f$: error
medio en el último segundo del tramo. Fuente: \texttt{informe/scripts/preparar\_datos.py} y
\texttt{etapaD/resultados/D1\_metricas.csv}.}
\label{tab:exp}
\centering
\footnotesize
\setlength{\tabcolsep}{3pt}
\begin{tabular}{@{}llrcrrr@{}}
\toprule
\textbf{Ensayo} & \textbf{Referencia} & $\Delta$ & $t_s$ & $M_p$ & $e_f$ & $\max|\thc|$ \\
 & (\si{\milli\metre}) & (\si{\milli\metre}) & (\si{\second}) & (\si{\percent}) & (\si{\milli\metre}) & (\si{\degree}) \\
\midrule
Prototipo, X & $0 \rightarrow 40$     & 40,7  & $> 4{,}7$ & 0 & 23,2 & 2,9 \\
Prototipo, X & $40 \rightarrow -40$   & -57,1 & $> 4{,}8$ & 0 & 27,8 & 1,9 \\
Prototipo, X & $-40 \rightarrow 0$    & 12,4  & $> 4{,}8$ & 0 & 9,3  & 1,3 \\
Prototipo, Y & $0 \rightarrow 40$     & 83,5  & $> 4{,}8$ & 0 & 80,5 & 2,0 \\
Prototipo, Y & $40 \rightarrow -40$   & 0,5   & $> 4{,}8$ & 0 & 38,6 & 1,6 \\
Prototipo, Y & $-40 \rightarrow 0$    & 76,7  & $> 4{,}8$ & 0 & 31,8 & 2,0 \\
\midrule
MIL D4, X & $0 \rightarrow 50$        & 51,3  & 4,00 & 14,1 & 1,5 & 2,1 \\
MIL D4, Y & $0 \rightarrow 50$        & 48,5  & 4,20 & 16,7 & 3,0 & 1,9 \\
\bottomrule
\end{tabular}
\end{table}

En el prototipo, ningún escalón se estableció dentro del tramo, así que $t_s$ solo tiene una
cota inferior, y el error al final fue de \SIrange{9}{80}{\milli\metre}, contra
\SIrange{1.5}{3.0}{\milli\metre} en el MIL. No hubo sobreimpulso: la bola se acerca lentamente y
se detiene antes de llegar. El eje Y fue el peor. En el escalón de 0 a \SI{40}{\milli\metre} la
bola estaba en \SI{-43}{\milli\metre} y prácticamente no se movió, aunque el plato llegó a pedir
\SI{2}{\degree}. Es el comportamiento que se espera de la zona muerta y la baja ganancia efectiva
identificadas en la Sección~II: con $\Kp = \SI{0.034}{\degree\per\milli\metre}$, un error de
\SI{30}{\milli\metre} pide apenas \SI{1}{\degree}, del orden de la zona muerta.

En el círculo, la órbita quedó centrada en $(6, -4)\,\si{\milli\metre}$, con un radio medio de
\SI{45}{\milli\metre} (percentiles 10 y 90 de \SI{16}{} y \SI{73}{\milli\metre}) y un error
mediano respecto de la referencia de \SI{30}{\milli\metre}. La bola estuvo a más de
\SI{115}{\milli\metre} del centro solo el \SI{0.2}{\percent} del tiempo, y se detectó en todos
los cuadros.

El 2026-11-18 se repitió el modo de centrado y apareció el mismo problema de forma más clara
(\texttt{sesion\_20261118\_092357}): la bola quedó \SI{55}{\second} quieta en
$(110, 58)\,\si{\milli\metre}$ con el plato pidiendo \SI{3.6}{\degree} y \SI{1.8}{\degree}, y el
integrador en cero, porque el firmware lo descargaba con la bola lejos de la referencia. Con eso
se corrigieron dos cosas: el \emph{trim} pasó a $\{1{,}5;\ 0{,}4\}\,\si{\degree}$, estimado como
la inclinación aplicada menos la zona muerta, y el integrador volvió a la regla de
\eqref{eq:ley}, que suma con la bola quieta en cualquier posición. En una simulación de un eje
con zona muerta, ese cambio dejó la bola a menos de \SI{13}{\milli\metre} del centro en los 30
casos probados, contra \SIrange{51}{57}{\milli\metre} en 3 casos con la regla anterior. Esta
última versión no se alcanzó a probar en el prototipo.

Con estos resultados, los requisitos R1 y R3 no se cumplen en la planta real. R4 tampoco se
cumplió siempre: en el log del 2026-11-18 la bola llegó al borde en $x \approx \SI{127}{\milli\metre}$,
y en las pruebas con la bola de PLA llegó a él decenas de veces. Fueron los bordes del plato los
que evitaron que cayera.
